\subsection{The source theta functions separate points and tangents}
\label{subsec:ftb-torus-theta-separation}

This starts the converse complex-torus construction.
It is an analytic argument and assumes no algebraic model of
the torus.  Let \(\tau=\tau^t\) with
\(Y=\operatorname{Im}\tau>0\), set
\(\Lambda_\tau=\mathbb Z^g+\tau\mathbb Z^g\), and write
\(e_C(w)=\exp(2\pi iw)\).
Let \(\mathcal T_N(\tau)\) be the entire functions with law
\[
 f(z+\tau n+m)=
 e_C\!\left(-N\{n^t\tau n/2+n^tz\}\right)f(z)
                       \quad(m,n\in\mathbb Z^g).         \tag{TS1}
\]
For \(N\ge1\) its classical basis is
\[
 T_{N,c}(z)=
 \theta_{c/N,0}(Nz,N\tau),
                     \qquad c\in\{0,\ldots,N-1\}^g .
 \tag{TS2}
\]
We use the already printed Gaussian convergence,
Fourier coefficient independence and coefficient-detection
argument: the \(N^g\) functions in (TS2) satisfy (TS1),
are linearly independent, and span \(\mathcal T_N(\tau)\).
Here no algebraic section-rank law is used.
Indeed a theta function, after its usual exponential
normalization, is integer-periodic; its Fourier coefficients
are determined by their \(N^g\) residue classes by the
\(\tau\)-quasi-period equations.
If those coefficients vanish, the Fourier uniqueness
argument makes the function zero.
This bounds the dimension by \(N^g\), and the \(N^g\)
independent functions (TS2) give the spanning assertion.

Put \(t(z)=\theta_{0,0}(z,\tau)\).
The case \(N=1\) shows that \(t\) is entire, nonzero, and
satisfies (TS1) at level one.
A nonzero entire function on a finite-dimensional complex
vector space has dense nonzero locus.
To see this in the form used below, vanishing on an open
ball would imply vanishing on each complex line from a
point in that ball to an arbitrary point, by the
one-variable identity theorem.
Its nonzero locus is also open, so finitely many such
nonzero loci have nonempty intersection.

\begin{lemma}[No base points from level two onward]
\label{lem:ftb-torus-theta-no-base-points}
For \(N\ge2\) and each \(z\in\mathbb C^g\), some
\(f\in\mathcal T_N(\tau)\) has \(f(z)\ne0\).
Consequently the basis in (TS2) has no common zero.
\end{lemma}
\begin{proof}
Write \(N=M+1\) with \(M\ge1\).
The two entire functions of \(u\),
\(t(z+u)\) and \(t(z-Mu)\), are nonzero functions:
translations and multiplication by nonzero \(M\) are onto.
Choose \(u\) in their common nonzero locus.
Then
\[
                   f(w)=t(w+u)^M t(w-Mu)
\]
is nonzero at \(z\).
The quasi-period laws give the level \(M+1=N\) law:
the shift contributions \(Mu\) and \(-Mu\) cancel in the
exponent, leaving
\(-N(n^t\tau n/2+n^tw)\).
Thus \(f\in\mathcal T_N(\tau)\).
Spanning by (TS2) now gives the last assertion.
\end{proof}

We recall the analytic translate-unit test, proved in the
course of Lemma~\ref{lem:ftb-actual-half-characteristic}.
Its proof is independent of the algebraic fibre in that
lemma:
\[
 \begin{gathered}
  U\text{ an entire unit},\qquad
  t(z+w)=U(z)t(z)\text{ for all }z\\
                 \Longrightarrow\quad w\in\Lambda_\tau .
 \end{gathered}                                           \tag{TS3}
\]
More explicitly, cancellation on the dense nonzero locus
of \(t\) gives
\(U(z+\tau n+m)=e_C(-n^tw)U(z)\).
The logarithmic derivative \(U^{-1}dU\) is entire and
\(\Lambda_\tau\)-periodic.
The lattice spans the real vector space, so the
fundamental-parallelepiped bound and Liouville make it
constant.  Thus \(U(z)=U(0)\exp(\ell(z))\).
Integer periods give
\(\ell(e_i)=2\pi i k_i,\ k_i\in\mathbb Z\).
The \(\tau e_i\) periods then give
\(w+\tau k\in\mathbb Z^g\), proving (TS3).
This records the actual translate test rather than an
assumption about a theta divisor stabilizer.

\begin{proposition}[Point separation]
\label{prop:ftb-torus-theta-point-separation}
Let \(N\ge3\).
If \(z,z'\in\mathbb C^g\) and \(\gamma\in\mathbb C\)
satisfy
\[
             f(z)=\gamma f(z')
                         \quad\text{for all }f\in\mathcal T_N(\tau),
                                                               \tag{TS4}
\]
then \(z-z'\in\Lambda_\tau\).
\end{proposition}
\begin{proof}
Write \(r=N-3\).
For all \(a,b,d\), the function
\[
 F_{a,b,d}(w)=
 t(w+a)t(w+b)t(w-(rd+a+b))t(w+d)^r
\]
has level \(N\), since the sum of its shifts, with their
multiplicities, is zero.
Equation (TS4) therefore gives
\[
 \begin{split}
 &t(z+a)t(z+b)t(z-(rd+a+b))t(z+d)^r\\
 &\hspace{12mm}=
 \gamma t(z'+a)t(z'+b)t(z'-(rd+a+b))t(z'+d)^r .
 \end{split}                                             \tag{TS5}
\]
Choose \(d\) with \(t(z+d)t(z'+d)\ne0\), using two dense
nonzero loci.
Choose \(a\) with \(t(z+a)\ne0\), and then \(b\) with both
remaining factors on the left nonzero.
This proves \(\gamma\ne0\), including \(r=0\).

Set \(\Delta=z-z'\).  For fixed \(b\) define entire functions
of \(x\) by
\[
 \begin{split}
 D_b(x)&=t(z+b)t(z-(rd+x-z'+b))t(z+d)^r,\\
 A_b(x)&=\gamma t(z'+b)t(z'-(rd+x-z'+b))t(z'+d)^r.
 \end{split}
\]
Substitute \(a=x-z'\) into (TS5).  It gives
\[
                    t(x+\Delta)D_b(x)=t(x)A_b(x).        \tag{TS6}
\]
For each \(x\), a choice of \(b\) makes both \(D_b(x)\)
and \(A_b(x)\) nonzero.
In detail, the four translated or reflected functions
of \(b\) occurring here are all nonzero entire functions;
their open dense nonzero sets have nonempty intersection.
For any \(b,b'\), (TS6) implies
\[
                       A_bD_{b'}=A_{b'}D_b
\]
on the dense nonzero set of \(t\), and then everywhere
by continuity.
Hence the locally holomorphic ratios \(A_b/D_b\), using
a \(b\) with nonzero denominator at each point, agree
on their overlaps.
Their numerators can be chosen nonzero at that point
too.  They glue to an entire unit \(U\) with
\(t(x+\Delta)=U(x)t(x)\).
The analytic test (TS3) gives \(\Delta\in\Lambda_\tau\).
\end{proof}

\begin{proposition}[Tangent separation]
\label{prop:ftb-torus-theta-tangent-separation}
Let \(N\ge3\).
If \(z,\xi\in\mathbb C^g\) and \(\gamma\in\mathbb C\)
satisfy
\[
       d f_z(\xi)=\gamma f(z)
                         \quad\text{for all }f\in\mathcal T_N(\tau),
                                                               \tag{TS7}
\]
then \(\xi=0\).
\end{proposition}
\begin{proof}
Write \(N=m+3\), \(k=m+1\), \(s=Nz\), and
\(D(a)=d t_a(\xi)\).
The derivative \(D\) is entire.
For arbitrary \(a,b\) set \(c=s-a-kb\).
Apply (TS7) to
\[
 x\longmapsto
 t(x+a-z)t(x+b-z)^k t(x+c-z),
\]
whose shifts sum to zero, so it has level \(N\).
The product rule gives the exact three-term identity
\[
 \begin{split}
 D(a)t(b)^kt(c)
 &+kt(a)t(b)^{k-1}D(b)t(c)
       +t(a)t(b)^kD(c)\\
                         &=\gamma t(a)t(b)^kt(c).        \tag{TS8}
 \end{split}
\]
We first construct an entire \(\psi\) with \(D=\psi t\).
Near any chosen \(a_0\), choose \(b\) such that
\(t(b)t(s-a_0-kb)\ne0\); multiplication by \(k\ne0\)
makes the second translated function nonzero, and dense
nonvanishing provides the common choice.
Near \(a_0\), (TS8) can then be solved as
\[
 D(a)=
 \left(\gamma-k\frac{D(b)}{t(b)}
              -\frac{D(s-a-kb)}{t(s-a-kb)}\right)t(a).
 \tag{TS9}
\]
The right coefficient is holomorphic there.
Two such local coefficients agree where \(t\ne0\),
hence on their overlap by continuity.
They glue to the required entire \(\psi\).

Substitute \(D=\psi t\) into (TS8), and cancel first on
the common nonzero locus of its three theta factors.
That locus is dense.  Continuity, first in \(a\) and then
in \(b\), gives
\[
                 \psi(a)+k\psi(b)+\psi(s-a-kb)=\gamma
                                             \quad(a,b\in\mathbb C^g).
 \tag{TS10}
\]
Compare (TS10) at
\((a,b)=(y+w,k^{-1}(s-y-w))\) and
\((y,k^{-1}(s-y-w))\).
The last arguments are respectively \(0,w\); subtraction
yields
\[
                  \psi(y+w)=\psi(y)+\psi(w)-\psi(0).      \tag{TS11}
\]

The source now integrates on each complex line, keeping
track of the entire unit thereby produced.
Choose an entire primitive \(K\) of
\(r\mapsto\psi(r\xi)-\psi(0)\), which exists on \(\mathbb C\).
The derivative of
\[
 r\longmapsto
       t(y+r\xi)\exp\{-r\psi(y)-K(r)\}
\]
is zero, using \(D=\psi t\) and (TS11).
Evaluating at \(r=0\) and \(r=u\) proves
\[
 t(y+u\xi)=
    \exp\{u\psi(y)+K(u)-K(0)\}\,t(y).                    \tag{TS12}
\]
For every \(u\in\mathbb C\), the exponential in (TS12)
is an entire unit as a function of \(y\).
By (TS3), \(u\xi\in\Lambda_\tau\) for every \(u\).
The map \(u\mapsto u\xi\), valued in the discrete lattice
\(\Lambda_\tau\), is continuous.
Since \(\mathbb C\) is preconnected it is constant.
Its values at \(0\) and \(1\) give \(\xi=0\).
\end{proof}

The conclusions have the exact projective interpretation
used by the source.  For \(N\ge3\),
\[
        z\longmapsto[T_{N,c}(z)]_c
\]
is defined everywhere, is constant on the lattice cosets,
separates distinct cosets, and has injective differential
on the quotient.
For the last assertion choose one nonzero coordinate at
\(z\).  A tangent annihilating all its affine coordinate
ratios has \(dT_{N,c}(\xi)=\gamma T_{N,c}(z)\) for the same
\(\gamma\) and every \(c\), hence for every level-\(N\)
theta function by spanning.  Tangent separation gives
\(\xi=0\).
At this stage these are analytic assertions about the
specified theta functions; the algebraic projective model,
its exact zero locus and its group law are subsequent
constructions.

\paragraph{Frozen-source correspondence.}
With prefix
\texttt{src/Development/AbelianVarietyTheory/Proofs/ToricGeometry/},
the classical basis and spanning are
\texttt{ClassicalThetaLevelAndLinearIndependent.lean}
(\texttt{ddff4de8eb961822383e688c1ffcc27550412858}),
\texttt{CardLePowOfLinearIndependent.lean}
(\texttt{20085cbe9a9a42f684df9924f67d1e8fcc10a49f})
and \texttt{MemSpanThetaOfLevel.lean}
(\texttt{1301044dc8de1a3eca303cef3a0c9e70e35a9335}).
The basic theta law is
\texttt{RiemannThetaLevelOneNeZero.lean}
(\texttt{8b538b68f2936a221c065ac9f914742b164fc905});
base-point freeness is
\texttt{ExistsApplySubSmulNeZeroAndApplyAddNeZero.lean}
(\texttt{41e334a200dc2f4522939e76d9fe724de5fa0f0b}),
\srcid{exists_level_apply_ne_zero_of_two_le}.
The translate test is
\texttt{MemPeriodLatticeOfThetaTranslateEqUnitMul.lean}
(\texttt{eed4913226755b63075a8f4dd698a5aa11894c89}).
The two separation proofs are
\texttt{LevelThreeSeparatesPoints.lean}
(\texttt{c68668f77d245fb186e1b7d6998fb0b4bef35f05}),
\srcid{level_three_separates_points}, and
\texttt{LevelThreeSeparatesTangentVectors.lean}
(\texttt{06fd2923a41705ca764c21d1b7528ac7581dfc28}),
\srcid{exists_unit_of_three_term},
\srcid{level_three_separates_tangent_vectors}.
These are the first analytic producers used in
\texttt{TorusAlgebraicOfSectionsOfVeryAmple.lean}
(\texttt{649eb8bec5c5c88a5881803e36fd8f8648c8a57b}),
rather than an invocation of a converse algebraization
theorem.
